\documentclass[conference]{IEEEtran}
\IEEEoverridecommandlockouts
\usepackage{cite}
\usepackage{amsmath,amssymb,amsfonts}
\usepackage{algorithmic}
\usepackage{graphicx}
\usepackage{textcomp}
\usepackage{xcolor}
\usepackage{booktabs}

\graphicspath{{figures/}}

\def\BibTeX{{\rm B\kern-.05em{\sc i\kern-.025em b}\kern-.08em
    T\kern-.1667em\lower.7ex\hbox{E}\kern-.125emX}}

\begin{document}

\title{LLM-Orchestrated Kinematic Adaptation (LOKA): An Immune System for Dynamic Humanoid Locomotion}

\author{\IEEEauthorblockN{Vineet Kulkarni}
\IEEEauthorblockA{\textit{School of Interactive Computing} \\
\textit{Georgia Institute of Technology}\\
Atlanta, Georgia, USA \\
vkulkarni46@gatech.edu}
\and
\IEEEauthorblockN{Nitish Sontakke}
\IEEEauthorblockA{\textit{School of Interactive Computing} \\
\textit{Georgia Institute of Technology}\\
Atlanta, Georgia, USA \\
nitishsontakke@gatech.edu}
\and
\IEEEauthorblockN{Dr.\ Sehoon Ha (Advisor)}
\IEEEauthorblockA{\textit{School of Interactive Computing} \\
\textit{Georgia Institute of Technology}\\
Atlanta, Georgia, USA \\
sehoonha@gatech.edu}
}

\maketitle

\begin{abstract}
Foundation models reason fluently about missions and failures; humanoid balance does not tolerate their latency.
Domain randomization and offline skill discovery build broad policies, yet struggle to isolate discrete mid-mission faults or accept operator language as a change to objectives and success criteria.
Code-as-policy LLM loops are too slow---and brittle---for underactuated bipeds.
We present \textbf{LOKA} (LLM-Orchestrated Kinematic Adaptation), a dual-rate \emph{immune architecture} for humanoid locomotion: a certified fast plant (centroidal MPC, whole-body quadratic control, and classical reflexes) owns millisecond survival, while a language model acts as a slow, auditable adapter that rewrites the plant's objective landscape, gait policy, physical self-belief, and mission criteria through typed interfaces.
Streaming proprioception is compressed into semantic tags; multi-turn episodes diagnose unmodeled change without privileged fault labels; interventions are improvement-gated and capped so stable residuals become new operating points rather than endless retuning.
LOKA does not invent footholds or compete with hierarchical push-recovery MPC on impulse magnitude; it targets \emph{mission continuity under discrete change}---payload, friction, actuator loss, and operator intent---where offline skill libraries and free-form code generation are poorly matched.
The same immune loop is framed as a plant-agnostic framework via typed adapters, with a Unitree~G1 stand/walk stack as the reference embodiment.
\end{abstract}

\begin{IEEEkeywords}
Humanoid Locomotion, Model Predictive Control, Large Language Models, Online Adaptation, Fault Recovery, Immune Architecture
\end{IEEEkeywords}

\section{Introduction}
\label{sec:introduction}

Humanoid robots are dynamically unstable and underactuated: balance lives on a timescale of tens of milliseconds, while foundation-model inference lives on a timescale of seconds.
Safely coupling language-level reasoning to locomotion therefore cannot mean putting an LLM in the torque loop.
Yet discrete, unmodeled events---a payload shift, an ice patch, a seized actuator---and sparse operator requests (``crouch and limp left'') routinely exceed what a fixed MPC cost schedule or an offline Domain-Randomization (DR) policy was prepared for~\cite{TobinDR}.

Existing paradigms leave a characteristic gap.
Vision-Language-Action (VLA) models map perception to actions but struggle with compositional generalization and long-horizon compounding error~\cite{OpenVLA, ACT_VLA}.
Unsupervised and language-guided \emph{skill discovery} (DIAYN, LSD, LGSD) builds diverse behavioral repertoires offline~\cite{DIAYN, LSD, LGSD}; Eureka-style systems search over reward \emph{code} before reinforcement learning~\cite{Eureka, Lang2Reward}.
These lines excel at repertoire learning, not at mid-flight surgical retuning of a certified balancer.
Hierarchical and capture-point MPC stacks push recovery impulse and step timing to the state of the art~\cite{HierarchicalPushMPC, ALIPWalkMPC}, but they do not natively interpret operator language or rewrite mission success criteria when the body itself changes.
Code-as-policy LLM loops~\cite{Lang2Reward, Eureka} are expressive yet latently and syntactically unsafe for bipeds: a single invalid program in the catch window is catastrophic~\cite{CodeAsPolicyBrittleness}.

\paragraph{Positioning.}
We introduce \textbf{LOKA}---LLM-Orchestrated Kinematic Adaptation---as an \emph{online immune system} for certified humanoid locomotion controllers.
The fast plant (reduced-order MPC, whole-body QP, and optional classical reflexes) never waits on the language model.
The LLM is a slow system-identification and mission layer: it compresses proprioceptive anomalies into semantic tags and emits a typed YAML scratchpad that mutates allowlisted impedance and friction knobs, task and gait setpoints, \texttt{Error\_Tracking} success criteria, and a private MuJoCo \emph{belief} model distinct from the hidden plant.
Millisecond recovery remains classical; LOKA is evaluated on post-event adaptation quality, operator mission retargeting, and cessation of unproductive retuning when residuals plateau.
Relative to skill discovery, LOKA does not learn latent skills \(\pi(a{\mid}s,z)\); relative to hierarchical MPC, it does not claim larger recoverable impulses; relative to code-as-policy, it forbids free-form programs in the balance loop.

The key contributions are:
\begin{itemize}
\item \textbf{Dual-rate immune contract:} A hardware-capable architecture that isolates high-frequency physical execution from LLM inference, with an explicit rule that anything required in ${<}300\,\mathrm{ms}$ cannot wait on language.
\item \textbf{Typed plant$\neq$belief adaptation:} A YAML scratchpad that mutates costs, gait policy, tasks, mission criteria, and kinematic belief---never code---through hard allowlists and affordance clamps suited to underactuated bipeds.
\item \textbf{Semantic telemetry compression:} Effort, contact, force-tracking, and mission excess compressed into tags that trigger multi-turn failure episodes without privileged fault labels.
\item \textbf{Residual intelligence:} Intervention caps and plateau detection that accept a stable mismatch as a new operating point (lean, belief, re-baselined mission nominal) rather than thrashing gains.
\item \textbf{Framework shape:} A plant-agnostic orchestrator intended to harden \emph{any} locomotion stack exposing torque control, a typed parameter surface, and telemetry---with a Unitree~G1 stand/walk reference implementation.
\end{itemize}

\section{Related Work and Literature Review}
\label{sec:related}

To position LOKA within the current robotics landscape, we analyze related architectures across four primary research domains.

\subsection{Slow-Fast Orchestration and Online Fault Recovery}
\label{sec:related-slowfast}

Decoupling symbolic cognitive reasoning from reactive physical stabilization has been explored across multiple domains.
The AESOP framework utilizes a high-frequency anomaly classifier to trigger a slow generative LLM reasoner, relying on an MPC to preserve physical feasibility during the reasoning latency gap on embedded hardware~\cite{AESOP}.
The REFLECT framework processes multisensory log data into a hierarchical natural language summary to enable an LLM to explain manipulation failures and guide corrective behavior~\cite{REFLECT}.
While structurally related, REFLECT operates as a post-mortem diagnostic tool for general tasks, whereas LOKA's Telemetry Compressor monitors high-frequency proprioceptive stream inputs to execute online adjustments.

Online parameter modification has been explored in wheeled robotics and autonomous driving.
For instance, the Hey Robot! framework uses an LLM to modulate objective cost weights based on human preferences for mobile navigation, though it lacks structural awareness of the physical model~\cite{HeyRobot}.
Similarly, RobotxR1 exposes MPC cost structures and safety margins as online-tunable parameters that a closed-loop reinforcement learning agent alters via language prompts~\cite{RobotxR1}.
However, RobotxR1 treats the underlying kinematic model as an unchangeable prior.
LOKA expands this authority by allowing structural modifications to the solver's body model.
The integration of neuro-symbolic planners, such as Recover, combines explicit rule-based logic with language agents to isolate physical faults~\cite{Recover}.
LOKA avoids the execution gaps of pure symbolic rule bases by mapping outputs directly into a continuous optimal control loop.

Decoupling processing frequencies has also been applied to vehicle routing and monolithic transformers.
AsyncDriver implements an asynchronous multi-threaded system to separate high-latency driving reasoning from real-time path planning~\cite{AsyncDriver}.
In end-to-end foundation policies, Fast-in-Slow (FiS-VLA) builds a unified dual-system policy that embeds high-frequency execution token heads directly inside a larger, low-frequency VLM backbone~\cite{FiS-VLA}.
LOKA shares these multi-rate design patterns but rejects black-box end-to-end actions, ensuring low-level physical safety by anchoring the fast loop to a mathematically transparent MPC solver.

\subsection{Language-to-Rewards and Code Generation}
\label{sec:related-code}

The dominant paradigm for bridging foundation models with optimal control relies on generating executable Python code.
The Eureka and Text2Reward architectures utilize LLMs as automated programmers to search over and debug dense reward functions for reinforcement learning agents~\cite{Eureka, Text2Reward}.
While highly effective, these frameworks operate strictly as heavy offline synthesis engines requiring thousands of environment rollouts.

For online application, the ASPIRE framework distills successful debugging sessions into a reusable text-based skill library using synchronous programmatic synthesis for manipulation tasks~\cite{ASPIRE}.
Similarly, the Language to Rewards (L2R) framework translates natural language text descriptors into functional reward scripts optimized in a loop by a MuJoCo MPC~\cite{Lang2Reward}.
Video2Reward derives reward metrics from target video behaviors for legged robots~\cite{Video2Reward}.
While expressive, these systems are fundamentally bounded by the autoregressive code generation bottleneck, exposing the system to syntax failures and unrecoverable latency during runtime anomalies.
LOKA diverges sharply from this line of work by replacing code generation entirely with a strictly typed parameter and model mutation engine.

\subsection{LLM-Guided Feedback Control Loops}
\label{sec:related-feedback}

A parallel vector of research integrates language models into active feedback control architectures.
The LLMs-guided Adaptive Compensator utilizes an LLM to dynamically output an error-tracking compensator, verifying the loop via stability criteria on soft and humanoid links~\cite{LLMadaptive}.
At the joint gain level, the LLM-Controller framework dynamically modulates nonlinear feedback properties while relying on the underlying baseline tracking architecture to preserve local stability~\cite{LLMcontroller}.
Automated frameworks like ControlAgent and AgenticControl leverage multi-agent language systems to design classical proportional-integral-derivative (PID) layouts by reasoning through phase margins and settling times~\cite{ControlAgent, AgenticControl}.

To bridge optimization with neural safety, architectures such as Never Too Rigid to Reach enforce Lyapunov-constrained reinforcement learning priors to guide virtual model control on flexible arms~\cite{NeverRigid}.
Furthermore, the Meta-Control architecture synthesizes entirely new control topologies depending on changing operational skills~\cite{MetaControl}.
LOKA shares the philosophy of leveraging language priors to modulate control characteristics, but shifts the intervention point.
Rather than operating at the joint-level error feedback or gain layer, LOKA mutates the predictive horizon, task-space priorities, and kinematic assumptions of an online preview optimization engine.

\subsection{Alternative Robot-Language Interface Modalities}
\label{sec:related-interfaces}

Exposing intermediate control spaces to language models remains a highly active area of study.
The SayTap framework utilizes binary foot-contact matrices as a structured intermediate interface to map language commands to quadrupedal locomotion policies~\cite{SayTap}.
OmniVIC couples VLMs with retrieval-augmented generation to continually reconfigure the stiffness and damping properties of a Variable Impedance Controller (VIC), applying explicit force bounds for safety~\cite{OmniVIC}.
For manipulation, the IROSA framework maps user prompts to the operational parameters of Kernelized Movement Primitives (KMPs) through tool calls, completely isolating the language agent from direct trajectory generation~\cite{IROSA}.

To handle complex contact profiles, the CompliantVLA-adaptor and GenCHiP frameworks inject language-derived force compliance, variable impedance values, and compliance boundaries directly into manipulation matrices to ensure safety during physical interaction~\cite{CompliantVLA, GenCHiP}.
LOKA generalizes this mapping philosophy to the full-body locomotion of underactuated bipeds, utilizing a structured YAML configuration block to map semantic objectives directly into high-degree-of-freedom tracking priorities.

\section{Methodology: The LOKA Architecture}
\label{sec:method}

LOKA is an \emph{immune layer}: the certified locomotion plant runs continuously; the language model intervenes sparsely when proprioceptive semantics, mission criteria, or an operator request indicate that the current objectives, costs, or self-model are wrong.
The pipeline consists of four tightly coupled phases across a dual-rate architecture (Fig.~\ref{fig:architecture}).
The reference plant is a Unitree~G1 centroidal MPC + whole-body QP stack with stand and walk modes; the orchestrator API is plant-agnostic given a typed adapter.

\begin{figure}[t]
  \centering
  \fbox{\parbox{0.92\linewidth}{\centering\vspace{1.0em}
  \textbf{Fast loop:} plant (MPC+WBC $\pm$ gait/reflexes)\\[0.4em]
  Telemetry buffer $\rightarrow$ Compressor / Error\_Tracking / anomaly\\[0.4em]
  Anomaly, plateau gate, or operator request $\rightarrow$ LLM (slow)\\[0.4em]
  YAML scratchpad $\rightarrow$ costs / gait / tasks / belief\\[0.4em]
  Affordance clamp $\rightarrow$ apply without blocking torque
  \vspace{1.0em}}}
  \caption{LOKA dual-rate immune architecture.
  The high-frequency plant continues while a background LLM mutates allowlisted costs, gait policy, task goals, mission criteria, and internal-model belief from compressed telemetry.
  Millisecond recovery never waits on inference.}
  \label{fig:architecture}
\end{figure}

\subsection{Phase 1: Effort-Driven Telemetry Compression}
\label{sec:telemetry}

Rather than triggering adaptation only after tracking error has corrupted the bipedal state vector, LOKA monitors energy expenditure within the low-level planner and mission criteria exposed by \texttt{Error\_Tracking}.
The MPC continuously generates an optimal control input vector representing commanded actuator torques.
The Telemetry Compressor calculates the Effort Deviation Integral ($E_{dev}$) over a sliding temporal window of length $W$:
\begin{equation}
E_{dev} = \int_{t-W}^{t} \| u_{\text{actual}}(\tau) - u_{\text{nominal}}(\tau) \|^2 d\tau
\label{eq:effort}
\end{equation}
When $E_{dev}$ violates a safety boundary---or when composite mission excess exceeds a trigger threshold---the Compressor freezes the high-frequency telemetry buffer, delta-masks nominal state variables, and synthesizes the anomaly into semantic tags.
For example, a joint exhibiting maximum torque saturation but near-zero joint velocity is tagged as \texttt{[SATURATED\_SEIZED]}, while high torque coupled with erratic velocity is tagged as \texttt{[SATURATED\_FLAILING]}.
Early plant-health scoring (tilt, support margin, force mismatch, lateral CoM bias) can open a failure episode before mission criteria trip, while intentional crouch residuals that leave balance healthy do not.

In the current implementation, compressed reports are organized into optional sections:
(0)~tracked state relative to live \texttt{Error\_Tracking} terms and empirical MissionNominal;
(1)~balance / contact summaries;
(2)~motor-load asymmetry and saturation cues;
(3)~directive / mission context.
Sections can be toggled for ablations that isolate which observations the orchestrator requires.

\subsection{Phase 2: Asynchronous Dual-Rate Execution}
\label{sec:async}

To prevent LLM inference latency from freezing the controller, LOKA implements a clean asynchronous partitioning scheme:
\begin{enumerate}
\item \textbf{Fast control loop (main thread):} Samples the true state vector from the physics engine (with future expansions accommodating state estimators), invokes the MPC/WBC (and gait) step, maps actions to actuators, and updates rolling history queues.
\item \textbf{Slow cognitive loop (background thread):} When an anomaly is flagged, during a periodic heartbeat, or when an operator issues a natural-language request, a background worker constructs the semantic prompt---including prior interventions in the current failure episode---and dispatches it to the orchestrator.
\end{enumerate}
While the LLM executes its forward pass, the main thread continues physics steps unhindered, relying on the last known valid configuration to keep the robot stable.
Classical reflexes (slip, sat-disable, capture step), when present, also live on the fast path; LOKA may later \emph{retune} their thresholds but does not author torques in the catch window.

\subsection{Phase 3: The YAML Scratchpad}
\label{sec:scratchpad}

The interface between the slow and fast threads is a strictly typed shared repository: the YAML scratchpad.
The underlying MPC continuously minimizes a quadratic task-space cost over a preview horizon $N$:
\begin{equation}
J = \sum_{k=0}^{N-1} \left( \| x_k - x_{\text{ref}} \|_Q^2 + \| u_k \|_R^2 \right) + \| x_N - x_{\text{ref}} \|_P^2
\label{eq:mpc_obj}
\end{equation}
where $Q$ and $R$ represent diagonal tracking and effort penalization matrices.
The LLM mutates this optimization landscape by returning a formatted block with optional sections (omit keys that are unchanged):
\begin{itemize}
\item \texttt{Semantic\_State}: chain-of-thought hypothesis and analysis before parameter edits;
\item \texttt{Controller\_Targets}: allowlisted dotted paths (e.g., friction belief, base impedance)---robot-tuned QP costs remain locked;
\item \texttt{Task\_Targets}: mission and gait setpoints (\texttt{height}, \texttt{lean\_*}, \texttt{yaw}, \texttt{gait.mode}, \texttt{gait.speed}, \texttt{gait.heading}, cadence, swing height, \ldots);
\item \texttt{Error\_Tracking}: composite trigger threshold plus weighted terms over \texttt{qpos}/\texttt{qvel} indices with modes such as \texttt{below\_target}, \texttt{above\_target}, \texttt{abs\_above}, and \texttt{abs\_deviation};
\item \texttt{Model\_Mutations}: continuous structural edits to the controller's private prior---actuator \texttt{gear}, geom \texttt{friction}, and body \texttt{mass}---resolved against MJCF names, independent of the hidden physical plant.
\end{itemize}
Multi-turn failure episodes surface prior diagnoses so the model must revise its hypothesis rather than repeat unchanged mutations.
Operator requests may arrive outside failure and can rewrite task goals; when the mission changes, \texttt{Error\_Tracking} must be updated to match.
A plateau / intervention-cap gate stops thrashing when plant-health severity fails to improve, accepting the residual into lean, mass belief, and a re-baselined MissionNominal.

\subsection{Phase 4: Affordance Clamping and Execution}
\label{sec:apply}

Upon completion of inference, the resulting string is captured via a thread-safe queue, parsed, and handed to the main loop.
Before targets are applied, they pass through a deterministic affordance filter.
If the language model proposes a parameter that violates hard kinematic constraints (e.g., a torso height beyond maximum leg extension), the filter clamps the value to the nearest physical boundary.
Once validated, the main loop updates model properties and controller settings.
If a joint is mutated to a gear ratio of $0.0$ (virtual amputation), middleware zeros subsequent control signals for that actuator index; limp recovery further requires WBC torque-authority and gait/task edits (and ideally a fast sat-disable reflex), not gear belief alone.
Mutations update what the controller \emph{believes}; they do not repair the physical plant.
Policy guidance requires aligning belief with observations---e.g., reducing gear on a dead motor---rather than increasing gear to ``force'' a seized joint.

\subsection{Locomotion plant: stand and walk}
\label{sec:plant}

The reference plant is selectable. The default, \texttt{legacy\_dcm}, couples a centroidal convex MPC (contact forces) with a whole-body QP (torques) and DCM capture-point footholds.
Three ALIP stacks (Xiong \& Ames, T-RO 2022) share the same \texttt{compute\_torque} contract: \texttt{alip\_footstep} keeps the MPC and WBC and replaces the foothold law; \texttt{alip\_paper} tracks height, torso and the swing foot with the WBC and does not treat horizontal CoM as an output; \texttt{alip\_wbqp} replaces both QPs with one contact-wrench program whose steps come from an ALIP MPC.
Walking adds a gait clock and a Cartesian swing-foot task tracked via foot-center Jacobians (quintic or Bézier lift; no separate analytical foot IK).
\texttt{gait.mode=walk} engages planned contact masks and velocity-aligned CoM references; \texttt{heading} and \texttt{yaw} provide travel direction and torso facing.
Foothold $xy$ and per-tick contact bits remain classical---out of the LLM vocabulary---so language commands such as ``walk at $0.3\,\mathrm{m/s}$ heading left'' retune policy knobs rather than invent geometry.

\subsection{Algorithm Overview}
\label{sec:algorithm}

The following pseudocode outlines the asynchronous orchestration loop.

\begin{algorithmic}[1]
\STATE \textbf{Initialize:} $Plant$, $LOKA_{memory}$, $Q_{thread}$
\STATE \textbf{Start Main Physics Loop} (high frequency)
\WHILE{Robot is Active}
    \STATE $E_{dev} \leftarrow$ Calculate\_Effort\_Deviation()
    \STATE $e_{track} \leftarrow$ Evaluate\_Error\_Tracking()
    \STATE $a_{early} \leftarrow$ Assess\_Plant\_Anomaly()
    \IF{($E_{dev} >$ Threshold \OR $e_{track} >$ Trigger \OR $a_{early}$ \OR OperatorRequest) \AND LLM not busy}
        \IF{PlateauOrCap(Episode)}
            \STATE Accept\_Residual(lean, belief, MissionNominal)
        \ELSE
            \STATE $Telemetry \leftarrow$ Compress\_State($Plant$)
            \STATE Spawn\_Thread(Query\_LLM($Telemetry$, EpisodeHistory), $Q_{thread}$)
        \ENDIF
    \ENDIF
    \IF{$Q_{thread}$ has message}
        \STATE $YAML_{updates} \leftarrow$ Parse\_Queue($Q_{thread}$)
        \STATE Validate\_Constraints($YAML_{updates}$)
        \STATE Apply\_Typed\_Mutations($YAML_{updates}$)
        \STATE Record\_Intervention($YAML_{updates}$)
    \ENDIF
    \STATE $u_{optimal} \leftarrow$ $Plant$.Step()
    \STATE Execute\_Torques(Mask\_Dead\_Actuators($u_{optimal}$))
\ENDWHILE
\end{algorithmic}

\section{Experimental Design and Evaluation}
\label{sec:experiments}

We evaluate LOKA in the MuJoCo physics simulation environment on an underactuated bipedal walker driven by MuJoCo MPC (MJPC), with planned extensions to additional humanoid embodiments.

\subsection{Evaluation Scenarios}
\label{sec:scenarios}

To validate robust reformulation across physical and semantic shifts, the system is subjected to four validation tasks:
\begin{itemize}
\item \textbf{Mass Shift:} An unmodeled mass of \(x\)\,kg (TODO) is attached asymmetrically mid-stride, forcing the system to detect effort imbalance and re-stabilize its center-of-mass target.
\item \textbf{Friction Shift:} The global surface friction coefficient is stepped down from a nominal value of $1.0$ to a reduced value (TODO) to evaluate real-time environmental adaptation.
\item \textbf{Hardware Failure:} A hard fault is injected into a critical actuator (e.g., right hip), setting physical torque capacity to zero to validate the virtual-amputation pipeline.
\item \textbf{Semantic / Operator Objectives:} The robot receives overarching language instructions (e.g., crouch and walk, crawl, hold upright) and/or must navigate a specialized corridor featuring a low overhanging ceiling followed by elevated rubble.
\end{itemize}

\subsection{Quantitative Performance Metrics}
\label{sec:metrics}

The framework is evaluated using the following metrics:
\begin{itemize}
\item \textbf{Task Success Rate (TSR):} Fraction of trials where the robot reaches the terminal criterion without falling.
\item \textbf{Time to Stabilization ($T_s$):} Duration between anomaly onset and restoration of a steady-state gait (or satisfaction of updated mission criteria).
\item \textbf{Effort Deviation Integral ($E_{dev}$):} Accumulated energy metric in Equation~\ref{eq:effort}, quantifying adaptation efficiency (alternative integrals over tracking residuals are also under consideration).
\item \textbf{Outer-loop turns:} Number of orchestrator interventions until success or timeout.
\item \textbf{Semantic Accuracy Percentage (SAP):} An automated judge (independent frontier model) comparing parsed scratchpad mutations against hidden fault metadata.
\item \textbf{Interface reliability:} Rates of invalid YAML and unresolved MJCF names.
\end{itemize}

\subsection{Comparative Baselines}
\label{sec:baselines}

\textbf{TODO:} Flesh out precise baseline implementations adapted to the dynamic bipedal domain (e.g., vanilla fixed MJPC, Domain Randomization without online mutation, ablations without model mutations or multi-turn history, and a code-generation / Language-to-Rewards variant where latency permits).

\begin{table}[t]
  \centering
  \caption{Placeholder results (fault recovery \%, operator success \%, median outer-loop turns).}
  \label{tab:results}
  \setlength{\tabcolsep}{3.5pt}
  \begin{tabular}{@{}lccc@{}}
    \toprule
    Method & Fault (\%) & Op.\ (\%) & Turns \\
    \midrule
    Fixed MJPC & -- & -- & -- \\
    Ablation (no mut.) & -- & -- & -- \\
    LOKA (full) & -- & -- & -- \\
    \bottomrule
  \end{tabular}
\end{table}

% \begin{figure}[t]
%   \centering
%   \includegraphics[width=\linewidth]{recovery_curves.pdf}
%   \caption{Tracking error vs.\ time after fault injection for LOKA and baselines.}
%   \label{fig:recovery}
% \end{figure}

\section{Discussion and Limitations}
\label{sec:discussion}

\paragraph{Interpretability.}
Hypotheses and analyses in the scratchpad provide a human-readable audit trail of recovery strategies, complementing numeric weight and model edits.

\paragraph{Latency and rate.}
LLM calls are orders of magnitude slower than MPC.
LOKA is therefore a supervisory layer for sparse events, not a replacement for the inner controller; asynchronous dual-rate execution is essential for stability during inference.

\paragraph{Safety and validity.}
Structured parsing, MJCF name resolution, and affordance clamping reduce free-form action risk, but incorrect mutations can still destabilize the plant.
Harder verification, sandboxing, and human approval remain important for hardware deployment.

\paragraph{Limitations.}
Results will depend on prompt quality and the chosen LLM; hallucinations of invalid parameter names can still occur.
The mutation vocabulary is intentionally small (gear, friction, mass).
Current evaluation focuses on a walker embodiment in simulation, with humanoid transfer and richer spatial sensing still in progress.
Some related citations remain placeholders pending final bibliographic cleanup.

\section{Conclusion}
\label{sec:conclusion}

LOKA couples an LLM semantic orchestrator with MuJoCo MPC to adapt costs, planner settings, task goals, mission criteria, and internal-model belief under faults and operator intent.
By compressing telemetry, isolating inference on a slow thread, and requiring multi-turn hypothesis revision through a typed YAML interface, LOKA aims to make adaptive locomotion both interpretable and operator-addressable---without the latency and brittleness of online code generation.
Future work includes completing the evaluation matrix, multi-embodiment tasks, stronger verification of proposed mutations, closed-loop hardware trials, and richer dynamics-parameter identification within the same orchestration loop.


\section*{Appendix}
\label{sec:appendix}

\subsection{System Prompt Formulation}
TODO: Insert the baseline system prompt instructing the language model on the structure of the YAML scratchpad, available MPC tuning parameters, error-tracking policy, and model-mutation constraints.

\subsection{Orchestrator Output Examples}
TODO: Insert example YAML scratchpad blocks demonstrating zero-shot recovery solutions generated during the evaluation tasks (e.g., virtual amputation parameter modifications).


\begin{thebibliography}{00}
\bibitem{PhysionGap} Anonymous et al., ``Probing the physical dynamics of Vision-Language Models,'' \textit{arXiv preprint arXiv:2510.06251}, 2025.
\bibitem{TobinDR} J. Tobin et al., ``Domain randomization for transferring deep neural networks from simulation to the real world,'' in \textit{IROS}, 2017.
\bibitem{OpenVLA} M. Kim et al., ``OpenVLA: An Open-Source Vision-Language-Action Model,'' \textit{arXiv preprint arXiv:2406.09246}, 2024.
\bibitem{ACT_VLA} Anonymous et al., ``Unleashing More Actions via Action Compositional Training for VLA Models,'' \textit{arXiv preprint arXiv:2607.00351}, 2026.
\bibitem{MPC_Rigidity} Anonymous et al., ``Analysis of Model Predictive Control in Changing Kinematic Environments,'' \textit{Placeholder Citation}, 2025.
\bibitem{Lang2Reward} W. Yu et al., ``Language to Rewards for Robotic Skill Synthesis,'' in \textit{CoRL}, 2023.
\bibitem{CodeAsPolicyBrittleness} Anonymous et al., ``Evaluating the Brittleness of Code-as-Policy Frameworks in High-Frequency Domains,'' \textit{Placeholder Citation}, 2025.
\bibitem{AESOP} Anonymous et al., ``Real-Time Anomaly Detection and Reactive Planning with LLMs (AESOP),'' in \textit{RSS}, 2024.
\bibitem{REFLECT} Anonymous et al., ``REFLECT: Summarizing Robot Experiences for Failure Explanation and Correction,'' in \textit{CoRL}, 2023.
\bibitem{RobotxR1} Anonymous et al., ``RobotxR1: Embodied Robotic Intelligence via Closed-Loop RL on LLMs,'' \textit{arXiv preprint arXiv:2505.03238}, 2025.
\bibitem{Recover} Anonymous et al., ``Recover: A Neuro-Symbolic Framework for Failure Detection and Recovery,'' \textit{arXiv preprint arXiv:2404.00756}, 2024.
\bibitem{AsyncDriver} Anonymous et al., ``AsyncDriver: Asynchronous LLM-Enhanced Planner for Autonomous Driving,'' in \textit{ECCV}, 2024.
\bibitem{FiS-VLA} Anonymous et al., ``Fast-in-Slow (FiS-VLA): A Dual-System Foundation Model,'' \textit{arXiv preprint arXiv:2506.01953}, 2025.
\bibitem{Eureka} Y. J. Ma et al., ``Eureka: Human-Level Reward Design via Coding LLMs,'' in \textit{ICLR}, 2024.
\bibitem{Text2Reward} Anonymous et al., ``Text2Reward: Automated Dense Reward Function Generation for RL,'' in \textit{ICLR}, 2024.
\bibitem{Video2Reward} Anonymous et al., ``Video2Reward: Generating Reward Functions from Videos for Legged Robots,'' \textit{arXiv preprint arXiv:2412.05515}, 2024.
\bibitem{LLMadaptive} Anonymous et al., ``LLMs-guided Adaptive Compensator,'' \textit{arXiv preprint arXiv:2507.20509}, 2025.
\bibitem{LLMcontroller} Anonymous et al., ``LLM-Controller: Dynamic Robot Control Adaptation Using LLMs,'' \textit{Robotics and Autonomous Systems}, 2025.
\bibitem{ControlAgent} Anonymous et al., ``ControlAgent: Automating Control System Design via LLM Agents + Domain Expertise,'' in \textit{NeurIPS}, 2024.
\bibitem{AgenticControl} Anonymous et al., ``AgenticControl: An Automated Control Design Framework Using LLMs,'' \textit{arXiv preprint arXiv:2506.19160}, 2025.
\bibitem{NeverRigid} Anonymous et al., ``Never Too Rigid to Reach: Adaptive Virtual Model Control with LLM- and Lyapunov-Based RL,'' \textit{arXiv preprint arXiv:2510.22892}, 2025.
\bibitem{MetaControl} Anonymous et al., ``Meta-Control: Automatic Model-Based Control Synthesis for Heterogeneous Robot Skills,'' \textit{arXiv preprint arXiv:2405.11380}, 2024.
\bibitem{SayTap} G. Ji et al., ``SayTap: Language to Quadrupedal Locomotion,'' in \textit{CoRL}, 2023.
\bibitem{OmniVIC} Anonymous et al., ``OmniVIC: Self-Improving Variable Impedance Controller with VLM In-Context Learning,'' \textit{arXiv preprint arXiv:2510.17150}, 2025.
\bibitem{IROSA} Anonymous et al., ``IROSA: Interactive Robot Skill Adaptation using Natural Language,'' \textit{arXiv preprint arXiv:2603.03897}, 2026.
\bibitem{CompliantVLA} Anonymous et al., ``CompliantVLA-adaptor: VLM-Guided Variable Impedance Action,'' \textit{arXiv preprint arXiv:2601.15541}, 2026.
\bibitem{GenCHiP} Anonymous et al., ``GenCHiP: Generating Robot Policy Code for High-Precision Contact-Rich Manipulation,'' \textit{IROS}, 2024.
\bibitem{LanguageMPC} Anonymous et al., ``LanguageMPC: LLMs as Decision Makers for Autonomous Driving,'' \textit{arXiv preprint arXiv:2310.03026}, 2023.
\bibitem{HeyRobot} Anonymous et al., ``Hey Robot! Personalizing Robot Navigation through MPC with an LLM,'' in \textit{ICRA}, 2025.
\bibitem{ASPIRE} Anonymous et al., ``Agentic Skill Programming through Iterative Robot Exploration (ASPIRE),'' \textit{NVIDIA GEAR Lab}, 2026.
\bibitem{DrEureka} Y. J. Ma et al., ``DrEureka: Language Model Guided Sim-To-Real Transfer,'' in \textit{RSS}, 2024.
\bibitem{DIAYN} B. Eysenbach et al., ``Diversity is All You Need: Learning Skills without a Reward Function,'' in \textit{ICLR}, 2019.
\bibitem{LSD} S. Park et al., ``Lipschitz-Constrained Unsupervised Skill Discovery,'' in \textit{ICLR}, 2022.
\bibitem{CSD} S. Park et al., ``Controllability-Aware Unsupervised Skill Discovery,'' in \textit{ICML}, 2023.
\bibitem{LGSD} Anonymous et al., ``Language Guided Skill Discovery,'' in \textit{ICLR}, 2025.
\bibitem{HierarchicalPushMPC} Y. Chen and Q. Nguyen, ``Adapting Gait Frequency for Posture-Regulating Humanoid Push-Recovery via Hierarchical Model Predictive Control,'' in \textit{ICRA}, 2025.
\bibitem{ALIPWalkMPC} Anonymous et al., ``Hierarchical Reduced-Order Model Predictive Control for Robust Locomotion on Humanoid Robots,'' \textit{arXiv preprint arXiv:2509.04722}, 2025.
\bibitem{RLaugmentedMPC} Y. Chen and Q. Nguyen, ``Learning Agile Locomotion and Adaptive Behaviors via RL-augmented MPC,'' in \textit{ICRA}, 2024.
\bibitem{Being0} Anonymous et al., ``Being-0: A Humanoid Robotic Agent with Vision-Language Models and Modular Skills,'' \textit{arXiv preprint arXiv:2503.12533}, 2025.
\bibitem{UniHSI} Z. Xiao et al., ``Unified Human-Scene Interaction via Prompted Chain-of-Contacts,'' in \textit{ICLR}, 2024.
\end{thebibliography}


\end{document}
